\scanpage{16}

18) Supposons que $C,C'$ \deleted{\unsure{$\in$}} \inserted{$(C^2=C'^2=\unsure{y|\gamma|\eta})$} définissent des polarisations de \deleted{\unsure{$G$}} \inserted{$\mathcal M=\operatorname{Rep}(G)$}.

Quand celles-ci sont-elles \uline{égales} ? Il \uline{suffit} qu'on ait $C'=gCg^{-1}$, avec $g\in G(\mathbb R)$, car alors pour $\varphi$ forme invariante donnée, on aura

\[
\varphi^{C'}(x,y)=\varphi(x,gCg^{-1}y)=\varphi(g^{-1}x,Cg^{-1}y)=\varphi^C(g^{-1}x,g^{-1}y);
\]

Cette condition est-elle aussi \uline{nécessaire} ? \unsure{Oui|On}, \unsure{si on suppose} \inserted{$G'$ \unsure{connexe|connexes}} et on trouve même \unsure{$=$} la conjugaison par un élément de $G'(\mathbb R)^0$ !

% slanted
\margin{\deleted{\unsure{Quand}} \uline{NB} $G'$ connexe $\unsure{\Longleftrightarrow|\Longleftarrow}$ $\unsure{G|C}$ connexe (grâce à l'existence du $C$ \unsure{avec|ou} $C^2=y$)}

Tout d'abord, la conjugaison des gpes maximaux dans un gpe de Lie connexe montre que, \deleted{\unsure{les}} si $C,C'$ sont deux élémts de $G'$ qui sont hodgiens (sans avoir même $C^2=C'^2$ !) il existe $g\in G'(\mathbb R)^0$ tel que

\[
gCg^{-1}C'^{-1}\in Z'(\deleted{\unsure{\mathbb R}}\mathbb R)\qquad (\unsure{Z|Z'}=\text{Centre de }G')
\]

[ \uline{NB} $(G'/Z')(\mathbb R)^0\longleftarrow G'(\mathbb R)^0$ est surjectif ! ].

Donc on aura $C'=zgCg^{-1}$, $g\in G'(\mathbb R)^0$, $z\in Z'(\mathbb R)$. Si $C'^2=C^2$, on aura d'ailleurs $z^2=1$ i.e. $z\in {}_2Z'(\mathbb R)$.

\deleted{On \unsure{se}} On peut donc supposer $C'=z\deleted{\unsure{g}}C$, $z\in {}_2Z(\mathbb R)$, et on est ramené à prouver que si $z\ne1$, alors $C'=zC$ et $C$ donnent des polarisations diff.

Or si $z\ne1$, il y aura une repr. de $G$ dans un $V\ne0$, où $z$ opère par $-\mathrm{id}_V$. Si $\varphi$ est une forme $C$-polarisante, $\varphi^C>0$, or $\varphi^{C'}=\varphi^{zC}=-\varphi^C<0$, absurde !

Donc on trouve que l'ensemble des polarisations sur $\mathcal M=\operatorname{Rep}(G)$ qui sont définissables par un $C$ (i.e. qui sont ``hodgiennes'') forme un \uline{torseur} sous $ {}_2Z(\mathbb R)$ : Si \deleted{\unsure{$\varphi$ est une}} $P$ est une polarisation sur $\mathcal M$ et $z\in {}_2Z(\mathbb R)$, alors $z.P$ est la polarisation

% archive sheet number: 15
